\documentclass[11pt]{article}

\usepackage{amsmath,amsthm,amssymb,amsfonts,mathrsfs,bbm}
\usepackage[margin=1in]{geometry}

\newcommand{\op}{\operatorname}

\newcommand{\C}{\mathbb{C}}
\newcommand{\cj}{\overline}
\newcommand{\F}{\mathbb{F}}
\newcommand{\mc}{\mathcal}
\newcommand{\R}{\mathbb{R}}
\newcommand{\tb}{\textbf}
\newcommand{\tn}{\textnormal}

\pagestyle{empty}

\theoremstyle{definition}
\newtheorem*{prop}{Proposition}

\makeatletter
\renewcommand*\env@matrix[1][*\c@MaxMatrixCols c]{%
\hskip -\arraycolsep
\let\@ifnextchar\new@ifnextchar
\array{#1}}
\makeatother

\begin{document}

\section*{Exam 1 Review}

Exam 1 will take place on Friday, March 7th in class and it will cover the following topics: 
\begin{itemize}	
	\item Systems of linear equations, including equivalent systems, row operations, types of solution sets, and homogeneous systems
	
	\item Matrix and vector notations for systems, including echelon forms and nonsingular matrices
	
	\item Vector operations, including addition, scaling, linear combinations, spanning sets, and linear independence
	
	\item Orthogonality, including inner products, norms, orthogonal/orthonormal sets, and the Gram-Schmidt process
	
	\item Matrix operations, including addition, scaling, transpose, adjoints, matrix multiplication, and matrix inverses.
\end{itemize}

\noindent The following list of practice problems will be similar to those on the exam:

\begin{enumerate}
	\item Describe the set of solutions of the following system of linear equations:
	\begin{align*}
		2x - 3y + z & = 0 \\
		4 x - 6y + 2z & = 0 \\
		x - y & = 1
	\end{align*}
	
	\item Let $\tb{x} \in \R^n$. Prove that $(-1) \tb{x} = - \tb{x}$ using known facts of vector addition and scaling.
	
	\item Find the simplified form of
	\begin{equation*}
		\operatorname{span}\left( \begin{bmatrix} 1 \\ - 1 \\ 1 \end{bmatrix}, \; \begin{bmatrix} 2 \\ -3 \\ 3 \end{bmatrix}, \; \begin{bmatrix} 0 \\ 1 \\ - 1 \end{bmatrix} \right)
	\end{equation*}
	
	\item Find all $t \in \R$ such that 
	\begin{equation*}
		\begin{bmatrix} -2 -t \\ 3 \end{bmatrix} \; \tn{and} \; \begin{bmatrix} 9 \\ - 2- t \end{bmatrix}
	\end{equation*}
	are linearly independent.
	
	\item Let $\tb{x},\tb{y} \in \R^n$ be such that $(\tb{x} + \tb{y}) \cdot (\tb{x} - \tb{y}) = 0$. Prove that $\|\tb{x}\| = \|\tb{y}\|$.
        \\
            \textbf{Answer}
        \\
        Since $(\tb{x} + \tb{y}) \cdot (\tb{x} - \tb{y}) = 0$, $(\tb{x} + \tb{y})$ and $(\tb{x} - \tb{y})$ are orthogonal. Hence, based on the Pythagorean theorem, and by setting $\alpha_1, \alpha_2 = 1$, we have:
        \begin{align*}
		& \left\| (\tb{x} + \tb{y}) +  (\tb{x} - \tb{y}) \right\|^2 = \|(\tb{x} + \tb{y})\|^2 + \|(\tb{x} - \tb{y})\|^2\\
        \Rightarrow & \left\| 2\tb{x} \right\|^2 = \left\| \tb{x} \right\|^2 + \left\| \tb{y} \right\|^2 -2\tb{x} \cdot\tb{y} + \left\| \tb{x} \right\|^2 + \left\| \tb{y} \right\|^2 +2\tb{x} \cdot\tb{y} \\
        \Rightarrow & 4\left\|\tb{x} \right\|^2 = 2\left\| \tb{x} \right\|^2 + 2\left\| \tb{y} \right\|^2 
        \\
        \Rightarrow & \left\|\tb{x} \right\|^2 = \left\| \tb{y} \right\|^2 
	\end{align*}
        Since $0 \leq \left\|\tb{x} \right\|, \left\| \tb{y} \right\|$:
        \[\left\|\tb{x} \right\| = \left\| \tb{y} \right\|\]
	
	\item  Use the Cauchy-Schwarz inequality to prove that if $0 < a,b,c,d$, then
	\begin{equation*}
		16 \leq (a + b + c + d) \left( \frac{1}{a} + \frac{1}{b} + \frac{1}{c} + \frac{1}{d} \right)
	\end{equation*}	
	
	\item Perform the Gram--Schmidt process on the following:
	\begin{equation*}
		\begin{bmatrix} 2 \\ 2 \\ 2 \end{bmatrix}, \begin{bmatrix} 1 \\ 0 \\ -1 \end{bmatrix}, \begin{bmatrix} 0 \\ 3 \\ 1 \end{bmatrix}
	\end{equation*}
        \\
            \textbf{Answer}
        \\
        Consider the given vectors in the following order:
        \begin{equation*}
		x_1 = \begin{bmatrix} 2 \\ 2 \\ 2 \end{bmatrix}, 
            x_2 = \begin{bmatrix} 1 \\ 0 \\ -1 \end{bmatrix}, 
            x_3 =\begin{bmatrix} 0 \\ 3 \\ 1 \end{bmatrix}
        \end{equation*}
        Now, suppose we want to construct the orthogonal set $v_1, v_2, v_3 \in \mathbb{R}^3$ using the Gram--Schmidt process. Based on this process, and proceeding inductively, we have:
        \[v_1 = x_1 = \begin{bmatrix} 2 \\ 2 \\ 2 \end{bmatrix}\]
        To calculate $v_2$:
        \begin{align*}
            v_2 
            & = x_2 - \frac{x_2 \cdot v_1}{v_1 \cdot v_1}v_1 \\
            & = \begin{bmatrix} 1 \\ 0 \\ -1 \end{bmatrix} - \frac{\begin{bmatrix} 1 \\ 0 \\ -1 \end{bmatrix} \cdot \begin{bmatrix} 2 \\ 2 \\ 2 \end{bmatrix}}{\begin{bmatrix} 2 \\ 2 \\ 2 \end{bmatrix} \cdot \begin{bmatrix} 2 \\ 2 \\ 2 \end{bmatrix}}\begin{bmatrix} 2 \\ 2 \\ 2 \end{bmatrix}\\
            & = \begin{bmatrix} 1 \\ 0 \\ -1 \end{bmatrix} - \frac{0}{12}\begin{bmatrix} 2 \\ 2 \\ 2 \end{bmatrix}\\
            & = \begin{bmatrix} 1 \\ 0 \\ -1 \end{bmatrix} 
        \end{align*}
        To calculate $v_3$:
        \begin{align*}
            v_3
            & = \begin{bmatrix} 0 \\ 3 \\ 1 \end{bmatrix} - \frac{x_3 \cdot v_1}{v_1 \cdot v_1}v_1 - \frac{x_3 \cdot v_2}{v_2 \cdot v_2}v_2 \\
            & = \begin{bmatrix} 0 \\ 3 \\ 1 \end{bmatrix} - \frac{\begin{bmatrix} 0 \\ 3 \\ 1 \end{bmatrix} \cdot \begin{bmatrix} 2 \\ 2 \\ 2 \end{bmatrix}}{\begin{bmatrix} 2 \\ 2 \\ 2 \end{bmatrix} \cdot \begin{bmatrix} 2 \\ 2 \\ 2 \end{bmatrix}}\begin{bmatrix} 2 \\ 2 \\ 2 \end{bmatrix} - \frac{\begin{bmatrix} 0 \\ 3 \\ 1 \end{bmatrix} \cdot \begin{bmatrix} 1 \\ 0 \\ -1 \end{bmatrix}}{\begin{bmatrix} 1 \\ 0 \\ -1 \end{bmatrix} \cdot \begin{bmatrix} 1 \\ 0 \\ -1 \end{bmatrix}}\begin{bmatrix} 1 \\ 0 \\ -1 \end{bmatrix}\\
            & = \begin{bmatrix} 0 \\ 3 \\ 1 \end{bmatrix} - \frac{8}{12}\begin{bmatrix} 2 \\ 2 \\ 2 \end{bmatrix} - \frac{-1}{2}\begin{bmatrix} 1 \\ 0 \\ -1 \end{bmatrix}\\ 
            & = \begin{bmatrix} 0 \\ 3 \\ 1 \end{bmatrix} - \begin{bmatrix} \frac{4}{3} \\ \frac{4}{3} \\ \frac{4}{3} \end{bmatrix} +\begin{bmatrix} \frac{1}{2} \\ 0 \\ -\frac{1}{2} \end{bmatrix}\\ 
            & = \begin{bmatrix} -\frac{5}{6} \\ \frac{5}{3} \\ -\frac{5}{6} \end{bmatrix}
        \end{align*}

		
	\item A complex matrix $\tb{A}$ is said to be \emph{self-adjoint} if $\tb{A}^\ast = \tb{A}$. 
	\begin{enumerate}
		\item Prove that 
		\begin{equation*}
			\frac{1}{2} \tb{A} + \frac{1}{2} \tb{A}^\ast
		\end{equation*}
		is self-adjoint
            \\
            \textbf{Answer}
            \\
            \begin{align*}
			(\frac{1}{2} \tb{A} + \frac{1}{2} \tb{A}^\ast)^\ast 
                & = (\frac{1}{2} \tb{A})^\ast + (\frac{1}{2} \tb{A}^\ast)^\ast \\
                & = (\cj{\frac{1}{2} \tb{A}})^\top + (\cj{\frac{1}{2} \tb{A}^\ast})^\top \\
                & = (\frac{1}{2} \cj{\tb{A}})^\top + (\frac{1}{2}\cj{ \tb{A}^\ast})^\top \\
                & = \frac{1}{2} (\cj{\tb{A}})^\top + \frac{1}{2}(\cj{ \tb{A}^\ast})^\top \\
                & = \frac{1}{2}\tb{A}^\ast + \frac{1}{2} (\tb{A}^\ast)^\ast \\
                & = \frac{1}{2}\tb{A}^\ast + \frac{1}{2} \tb{A} \\
                & = \frac{1}{2} \tb{A} + \frac{1}{2}\tb{A}^\ast
		\end{align*}
		
		\item Prove that
		\begin{equation*}
			\tb{A}^\ast \tb{A}
		\end{equation*}
		is self-adjoint
            \\
            \textbf{Answer}
            \\
            \begin{align*}
            (\tb{A}^\ast \tb{A})^\ast =\tb{A}^\ast(\tb{A}^\ast)^\ast =\tb{A}^\ast\tb{A}
            \end{align*}

	\end{enumerate}
	
	\item Solve the following system
	\begin{align*}
		7 x + 3y & = 1 \\
		5 x + 2y & = 1
	\end{align*}
	using matrix inverses.
        \\
        \textbf{Answer}
        \\
        As a matrix equation $\tb{A} \tb{x} = \tb{b}$ the following system can be written as:
	\begin{equation*}
		\begin{bmatrix} 7 & 3 \\ 5 & 2 \end{bmatrix} \begin{bmatrix} x \\ y \end{bmatrix} = \begin{bmatrix} 1 \\ 1 \end{bmatrix}
	\end{equation*}
        The inverse of $\tb{A}$ is given by
	\begin{equation*}
		\frac{1}{7\cdot2 - 5\cdot3} \begin{bmatrix} 2 & -3 \\ -5 & 7 \end{bmatrix} = -1\begin{bmatrix} 2 & -3 \\ -5 & 7 \end{bmatrix} = \begin{bmatrix} -2 & 3 \\ 5 & -7 \end{bmatrix}
	\end{equation*}
	and thus
	\begin{equation*}
		\begin{bmatrix} x \\ y \end{bmatrix} = \begin{bmatrix} -2 & 3 \\ 5 & -7 \end{bmatrix}  \begin{bmatrix} 1 \\ 1 \end{bmatrix} = \begin{bmatrix} 1 \\ -2 \end{bmatrix}
	\end{equation*}
	
	\item Find the inverse of the matrix
	\begin{equation*}
		\tb{A} = \begin{bmatrix} 0 & 1 & 2 \\ 1 & 0 & 3 \\ 4 & -3 & 8 \end{bmatrix}
	\end{equation*}
        \\
        \textbf{Answer}
        \\
        For any $1 \leq k \leq 3$ the columns of $\tb{A}^{-1}$ can be found by solving the system
	\begin{equation*}
		\tb{A} \tb{x} = \tb{e}_k
	\end{equation*}
	This can be done simultaneously as an augmented matrix as 
    follows: 
	\begin{equation*}
		\begin{bmatrix}[ccc|ccc] 0 & 1 & 2 & 1 & 0 & 0 \\ 1 & 0 & 3& 0 & 1 & 0 \\ 4 & -3 & 8 & 0 & 0 & 1 \end{bmatrix} \longrightarrow \begin{bmatrix}[ccc|ccc] 1 & 0 & 0 & -4.5 & 7 & -1.5 \\ 0 & 1 & 0 & -2 & 4 & -1 \\ 0 & 0 & 1 & 1.5 & -2 & 0.5 \end{bmatrix}
	\end{equation*}
    Hence
    \begin{equation*}
    \tb{A}^{-1} = \begin{bmatrix} -4.5 & 7 & -1.5 \\ -2 & 4 & -1 \\ 1.5 & -2 & 0.5 \end{bmatrix}
    \end{equation*}

    \item Show
    \[(\tb{A} \tb{B})^\ast = \tb{B}^\ast \tb{A}^\ast\]
    \\
        \textbf{Answer}
    \\
    \begin{align*}
		((\tb{A} \tb{B})^\ast)_{j,k}
            & = (\cj{(\tb{A} \tb{B})^{\top}})_{j,k} \\
            & = \cj{(\tb{A} \tb{B})^{\top}_{j,k}} \\
            & = \cj{(\tb{B}^{\top} \tb{A}^{\top} )_{j,k}} \\
            & = \cj{\sum_{h=1}^p\tb{B}^{\top}_{j,h}\cdot \tb{A}^{\top}_{h,k}} \\
            & = \sum_{h=1}^p\cj{\tb{B}^{\top}_{j,h}\cdot \tb{A}^{\top}_{h,k}} \\
            & = \sum_{h=1}^p\cj{\tb{B}^{\top}_{j,h}} \cdot \cj{\tb{A}^{\top}_{h,k}} \\ 
            & = \sum_{h=1}^p(\cj{\tb{B}^{\top}})_{j,h} \cdot (\cj{\tb{A}^{\top}})_{h,k} \\
            & = \sum_{h=1}^p(\tb{B}^\ast)_{j,h} \cdot (\tb{A}^\ast)_{h,k} \\
            & = \tb{B}^\ast \tb{A}^\ast
		\end{align*}

\end{enumerate}

\end{document}